\section*{Chapter \ref{ch:b2:palaeoenvironments} --- Fossils and Palaeoenvironments}

\begin{solution}{exo:b2:palaeoenvironments:1}
(a) $112/135 = 0.830$ and $112/150 = 0.747$. (b) The dead shells on the floor gather over years
or more: they include species that lived there earlier, or rarely, or were carried in, so the
\omterm{def:b2:palaeoenvironments:assemblages}{death assemblage} averages many generations of the community. (c) Soft-bodied groups, in
particular detritivores and carnivores such as worms: without a skeleton they decay before burial.
\end{solution}

\begin{solution}{exo:b2:palaeoenvironments:2}
(a) $5 \times 165 = \qty{825}{kyr}$, about \qty{0.8}{Myr}, within the 0.7 to 1 million years.
(b) $183.73 - 182.77 = \qty{0.96}{Myr}$, and $0.96/0.165 = 5.8$: about six horizons. (c) A
lineage changes within much less than a million years; only horizons separate the successive
forms, while a zone lumps them together.
\end{solution}

\begin{solution}{exo:b2:palaeoenvironments:3}
(a) Nereites \omterm{def:b2:palaeoenvironments:ichnofacies}{ichnofacies}: the deep-sea floor of a turbidite system, beyond the shelf break, where
food is scarce. (b) Skolithos \omterm{def:b2:palaeoenvironments:ichnofacies}{ichnofacies}: foreshore or shoreface. (c) Cruziana \omterm{def:b2:palaeoenvironments:ichnofacies}{ichnofacies}:
offshore, between the fair-weather and the storm wave bases. (d) Zoophycos \omterm{def:b2:palaeoenvironments:ichnofacies}{ichnofacies}: in the
Mesozoic, the lower shelf or the slope, below the storm wave base.
\end{solution}

\begin{solution}{exo:b2:palaeoenvironments:4}
(a) Bed A holds stenohaline groups: normal sea water. Bed B, poor in species but rich in
individuals, with plant debris and none of the stenohaline groups, suggests brackish water, for
example an estuary. (b) The missing groups are several, common and well preserved in marine beds
of the same age: their joint absence is a pattern, whereas one fossil can be missing by chance or
by poor preservation. (c) The sediment (tidal or river channels, mud drapes, roots), the \omterm{def:b2:palaeoenvironments:fossil-kinds}{trace
fossils} (a Cruziana \omterm{def:b2:palaeoenvironments:ichnofacies}{ichnofacies} of few species), and the passage of bed B, laterally, into
marine or into river deposits.
\end{solution}

\begin{solution}{exo:b2:palaeoenvironments:5}
(a) $P = 1/2$ when $kh/s = \ln 2$: $s = 0.05 \times 0.1/0.693 = \qty{0.0072}{m.yr^{-1}}$,
\qty{7.2}{mm} a year. (b) Halving $s$ doubles the exponent: $P = (1/2)^2 = 0.25$, half as much.
(c) $k$ is halved: $P = e^{-2.5} = 0.082$ instead of $e^{-5} = 0.0067$, $e^{2.5} = 12$ times more
likely. (d) Thick, robust shells are over-represented, and the more so the slower the burial.
\end{solution}

\begin{solution}{exo:b2:palaeoenvironments:6}
(a) The upper bound is $183.73 + 0.35 = \qty{184.08}{Ma}$: the chart's \qty{184.2}{Ma} lies
outside it, by \qty{0.12}{Myr}. (b) $183.73 - 182.77 = \qty{0.96}{Myr}$. (c) $6/0.96 =
\qty{6.25}{m}$ per million years. (d) A biozone is defined by fossils, whose appearances and
disappearances are close to simultaneous over wide areas; the thickness of rock laid down during
that time depends on the local rate of accumulation.
\end{solution}

\begin{solution}{exo:b2:palaeoenvironments:7}
(a) $0.49\,e^{-1046/3704} = 0.369$ and $0.49\,e^{-1166/3704} = 0.358$; mean $0.364$. (b) $120
\times (1 - 0.364) = \qty{76.3}{m}$. (c) $76.3/(1 - 0.49) = \qty{149.6}{m}$. (d) The factor
supposes the whole restored layer has the porosity $\varphi_0$ of the surface; but a layer
\qty{150}{m} thick at the surface is less porous at its base than at its top, so the same grains
fill less than \qty{149.6}{m}.
\end{solution}

\begin{solution}{exo:b2:palaeoenvironments:8}
(a) A unit volume holds $\varphi$ of water and $1 - \varphi$ of grains: $\rho = \varphi
\rho_{\text{water}} + (1 - \varphi)\rho_{\text{grain}}$, which gives the formula; $(2710 -
2573)/(2710 - 1000) = 0.080$. (b) $0.51\,e^{-956/4545} = 0.413$. (c) Cementation: calcite
precipitated in the pores, which carbonates undergo besides compaction, removes far more porosity
than burial alone.
\end{solution}

\begin{solution}{exo:b2:palaeoenvironments:9}
(a) Deep-sea turbidites (Nereites \omterm{def:b2:palaeoenvironments:ichnofacies}{ichnofacies}); a Jurassic floor of the lower shelf or slope
(Zoophycos); an offshore floor between the wave bases (Cruziana); a shoreface (Skolithos); a
swamp on land. (b) The water shallowed upward: the coast moved seaward. (c) A fall of sea level,
or a supply of sediment faster than the subsidence, as in front of a growing delta; a fall would
be found at the same time in other basins, a delta would be local and its sediment would come
from one side. (d) In a conformable succession neighbouring environments follow each other; a
jump over the intermediate ones points to a gap in time or to a fault.
\end{solution}

\begin{solution}{exo:b2:palaeoenvironments:10}
(a) The time to burial is $h/s$ = 100, 10 and \qty{1}{yr}. Soft body: $e^{-100} \approx
\num{4e-44}$, $e^{-10} = \num{4.5e-5}$, $e^{-1} = 0.37$. Shell ($k = \qty{0.01}{yr^{-1}}$):
$e^{-1} = 0.37$, $e^{-0.1} = 0.90$, $e^{-0.01} = 0.99$. (b) $s = kh/\ln 100 = 0.1/4.605 =
\qty{0.022}{m.yr^{-1}}$, \qty{22}{mm} a year. (c) Hardly, over months; but a storm bed, a
turbidite or an ash fall buries the floor in hours. (d) The constant rate of sedimentation: the
time to burial is then much shorter than $h/s$ with the mean rate, and $P = e^{-kT}$, with $T$
the actual time to burial, tends to one when $T$ is much less than $1/k$.
\end{solution}

\begin{solution}{exo:b2:palaeoenvironments:11}
(a) A thin layer $\mathrm dz$ at the depth $z$ holds the grains $(1 - \varphi(z))\,\mathrm dz$;
at the surface they fill $\mathrm dz_0$ with $(1 - \varphi_0)\,\mathrm dz_0 = (1 -
\varphi(z))\,\mathrm dz$. (b) $\varphi = 0.63\,e^{-816/1960} = 0.416$; $(1 - 0.416)/0.37 =
1.578$, so \qty{157.8}{m}. (c) \qty{9.2}{m} too much: as in \cref{exo:b2:palaeoenvironments:7},
the restored layer is not uniformly as porous as its surface. (d) The rate of grains,
$h(1-\varphi)/\Delta t$, uses only the present thickness and porosity, both measured.
\end{solution}

\begin{solution}{exo:b2:palaeoenvironments:12}
(a) $9.5/0.165 = 57.6$ horizons. (b) $180/57.6 = \qty{3.1}{m}$ of limestone today, $338/57.6 =
\qty{5.9}{m}$ of fresh sediment. (c) At most \qty{1.5}{m}: $180/1.5 = 120$ intervals, 121
samples. (d) Breaks in deposition, beds where ammonites are absent, and horizons condensed into
a few centimetres where the sediment arrived slowly.
\end{solution}

\begin{solution}{pb:b2:palaeoenvironments:1}
\textbf{1.} $866 \sin \ang{12} = 866 \times 0.2079 = \qty{180.0}{m}$.
\textbf{2.} $481 \times 0.2079 = \qty{100.0}{m}$.
\textbf{3.} The sandstone, by superposition: the limbs dip at only \ang{12}, so the beds are not
overturned and the lowest is the oldest.
\textbf{4.} $866 - 766 = \qty{100}{m}$ and $1046 - 866 = \qty{180}{m}$.
\textbf{5.} At least \qty{766}{m} there; its original top is eroded, so its full thickness is
unknown.
\textbf{6.} 0.70 at the surface; $0.70\,e^{-766/1408} = 0.406$ at \qty{766}{m}.
\textbf{7.} They are stenohaline groups: the water had the salinity of the open sea.
\textbf{8.} A Cruziana \omterm{def:b2:palaeoenvironments:ichnofacies}{ichnofacies}: the offshore, between the fair-weather and the storm wave
bases.
\textbf{9.} Storms: their waves reached the floor, broke the shells and laid them down in beds
with erosive bases.
\textbf{10.} The bottom water was poor in oxygen: few burrowers, laminae kept, pyrite formed.
\textbf{11.} The Zoophycos \omterm{def:b2:palaeoenvironments:ichnofacies}{ichnofacies}, in the Mesozoic the lower shelf or the slope, below the
storm wave base; enough oxygen at times for burrowers.
\textbf{12.} No: ammonites swam in the water above; their shells sank or drifted to the floor
after death and form a \omterm{def:b2:palaeoenvironments:assemblages}{death assemblage}.
\textbf{13.} $184.2 - 174.7 = \qty{9.5}{Myr}$.
\textbf{14.} $180/9.5 = \qty{18.9}{m}$ per million years.
\textbf{15.} $(2710 - 2573)/(2710 - 1000) = 0.080$.
\textbf{16.} $180 \times (1 - 0.080)/(1 - 0.51) = \qty{338}{m}$.
\textbf{17.} $338/9.5 = \qty{35.6}{m}$ per million years.
\textbf{18.} $180 \times 0.92/9.5 = \qty{17.4}{m}$ per million years.
\textbf{19.} The shale gave 11.3 and \qty{4.4}{m} per million years: the rate of grains fell
about four times. The carbonate of the shoal and the shelf stopped arriving as the floor
deepened, and only a slow rain of mud reached it.
\textbf{20.} It advanced: rivers, then a shoal, an offshore floor, a floor below the storm wave
base, then the open sea of the chalk. The column records a transgression.
\textbf{21.} $57 - 41 = \qty{16}{m}$ over the limestone, $86 - 57 = \qty{29}{m}$ over the shale,
\qty{45}{m} in all.
\textbf{22.} $(60 - 5) + 338 + 148.6 = \qty{542}{m}$, in $184.2 - 161.5 = \qty{22.7}{Myr}$:
\qty{23.9}{m} per million years.
\textbf{23.} $45/542 = 0.083$: about \qty{8}{\percent}, and less if the cycle chart exaggerates
the amplitude. The rest came from the subsidence of the floor.
\textbf{24.} The lime mud of the Varden limestone accumulated at about
\textbf{\qty[text-series-to-math=true]{36}{m}} per million years when it was laid down.
\end{solution}
